% Options for packages loaded elsewhere
\PassOptionsToPackage{unicode}{hyperref}
\PassOptionsToPackage{hyphens}{url}
\PassOptionsToPackage{dvipsnames,svgnames,x11names}{xcolor}
\documentclass[
  10pt,
  fontset=fandol]{ctexart}
\usepackage{xcolor}
\usepackage[margin=16mm]{geometry}
\usepackage{amsmath,amssymb}
\setcounter{secnumdepth}{-\maxdimen} % remove section numbering
\usepackage{iftex}
\ifPDFTeX
  \usepackage[T1]{fontenc}
  \usepackage[utf8]{inputenc}
  \usepackage{textcomp} % provide euro and other symbols
\else % if luatex or xetex
  \usepackage{unicode-math} % this also loads fontspec
  \defaultfontfeatures{Scale=MatchLowercase}
  \defaultfontfeatures[\rmfamily]{Ligatures=TeX,Scale=1}
\fi
\usepackage{lmodern}
\ifPDFTeX\else
  % xetex/luatex font selection
\fi
% Use upquote if available, for straight quotes in verbatim environments
\IfFileExists{upquote.sty}{\usepackage{upquote}}{}
\IfFileExists{microtype.sty}{% use microtype if available
  \usepackage[]{microtype}
  \UseMicrotypeSet[protrusion]{basicmath} % disable protrusion for tt fonts
}{}
\makeatletter
\@ifundefined{KOMAClassName}{% if non-KOMA class
  \IfFileExists{parskip.sty}{%
    \usepackage{parskip}
  }{% else
    \setlength{\parindent}{0pt}
    \setlength{\parskip}{6pt plus 2pt minus 1pt}}
}{% if KOMA class
  \KOMAoptions{parskip=half}}
\makeatother
\setlength{\emergencystretch}{3em} % prevent overfull lines
\providecommand{\tightlist}{%
  \setlength{\itemsep}{0pt}\setlength{\parskip}{0pt}}
\usepackage{amsmath,amssymb,mathtools,needspace}
\xeCJKDeclareCharClass{CJK}{"2032,"0400->"04FF,"2160->"216B,"2460->"2473,"25A0,"25C7,"25CB,"25CF}
\IfFontExistsTF{Arial Unicode MS}{\xeCJKsetup{AutoFallBack=true}\setCJKfallbackfamilyfont{\CJKrmdefault}{Arial Unicode MS}}{}
\setcounter{secnumdepth}{0}
\setlength{\emergencystretch}{2em}
\allowdisplaybreaks
\usepackage{bookmark}
\IfFileExists{xurl.sty}{\usepackage{xurl}}{} % add URL line breaks if available
\urlstyle{same}
\hypersetup{
  pdftitle={算式的建立},
  colorlinks=true,
  linkcolor={Maroon},
  filecolor={Maroon},
  citecolor={Blue},
  urlcolor={Blue},
  pdfcreator={LaTeX via pandoc}}

\title{算式的建立}
\author{}
\date{}

\begin{document}
\maketitle

\subsubsection{11.3.2　算式的建立}\label{ux7b97ux5f0fux7684ux5efaux7acb}

现在运用消元手续具体建立上述奇偶二分法的算式。回到递推问题（11.3.1），利用它的第 \(i-1\) 式从其第 \(i\) 式中消去 \(x_{i-1}\)，得

\[
\begin{cases}
x_i=b_i^{(1)},& i=0,1,\\
x_i=a_i^{(1)}x_{i-2}+b_i^{(1)},& i=2,3,\cdots,N-1,
\end{cases}
\tag{11.3.4}
\]

式中

\[
a_i^{(1)}=a_i a_{i-1},\quad i=2,3,\cdots,N-1,
\tag{11.3.5}
\]

\href{../../source-images/pdf-287.jpeg}{原页图像}

\[
b_i^{(1)}=
\begin{cases}
b_i,& i=0,\\
b_i+a_i b_{i-1},& i=1,2,\cdots,N-1.
\end{cases}
\tag{11.3.6}
\]

容易看出，式（11.3.4）可按下标的奇偶拆成两个规模减半的子问题，即

\[
\begin{cases}
x_0=b_0^{(1)},\\
x_{2i}=a_{2i}^{(1)}x_{2i-2}+b_{2i}^{(1)},\quad i=1,2,\cdots,N_1-1;
\end{cases}
\]

\[
\begin{cases}
x_1=b_1^{(1)},\\
x_{2i+1}=a_{2i+1}^{(1)}x_{2i-1}+b_{2i+1}^{(1)},\quad i=1,2,\cdots,N_1-1.
\end{cases}
\]

它们分别对应于形如式（11.3.2）的奇偶相关链。

前已指出，二分 \(k\) 步后加工得出的相关链式（11.3.3）含有 \(2^k\) 个结果 \(x_i\)，\(i=0,1,\cdots,2^k-1\)，且其步长增至 \(2^k\)，因此其相应的递推问题具有如下形式：

\[
\begin{cases}
x_i=b_i^{(k)},& i=0,1,\cdots,2^k-1,\\
x_i=a_i^{(k)}x_{i-2^k}+b_i^{(k)},& i=2^k,2^k+1,\cdots,N-1.
\end{cases}
\tag{11.3.7}
\]

为了导出系数 \(a_i^{(k)},b_i^{(k)}\) 的计算公式，回到前一步加工得出的递推问题：

\[
\begin{cases}
x_i=b_i^{(k-1)},& i=0,1,\cdots,2^{k-1}-1,\\
x_i=a_i^{(k-1)}x_{i-2^{k-1}}+b_i^{(k-1)},& i=2^{k-1},2^{k-1}+1,\cdots,N-1.
\end{cases}
\tag{11.3.8}
\]

显然，据此可得出 \(2^{k-1}\) 个新结果

\[
x_i=a_i^{(k-1)}b_{i-2^{k-1}}^{(k-1)}+b_i^{(k-1)},\quad i=2^{k-1},2^{k-1}+1,\cdots,2^k-1,
\]

与式（11.3.7）比较系数，有

\[
b_i^{(k)}=
\begin{cases}
b_i^{(k-1)},& i=0,1,\cdots,2^{k-1}-1,\\
b_i^{(k-1)}+a_i^{(k-1)}b_{i-2^{k-1}}^{(k-1)},& i=2^{k-1},2^{k-1}+1,\cdots,2^k-1.
\end{cases}
\tag{11.3.9}
\]

又将式（11.3.8）直接代入，得

\[
\begin{aligned}
x_i
&=a_i^{(k-1)}\left(a_{i-2^{k-1}}^{(k-1)}x_{i-2^k}+b_{i-2^{k-1}}^{(k-1)}\right)+b_i^{(k-1)}\\
&=\left(a_i^{(k-1)}a_{i-2^{k-1}}^{(k-1)}\right)x_{i-2^k}+\left(a_i^{(k-1)}b_{i-2^{k-1}}^{(k-1)}+b_i^{(k-1)}\right).
\end{aligned}
\]

再与式（11.3.7）比较系数，知

\[
\left\{
\begin{aligned}
a_i^{(k)}&=a_i^{(k-1)}a_{i-2^{k-1}}^{(k-1)},\\
b_i^{(k)}&=b_i^{(k-1)}+a_i^{(k-1)}b_{i-2^{k-1}}^{(k-1)},
\end{aligned}
\right.
\quad i=2^k,2^k+1,\cdots,N-1.
\]

将计算 \(b_i^{(k)}\) 的上述两组算式归并在一起，即可归纳出求解递推问题的奇偶二分算法 11.4。

\textbf{算法 11.4}　对 \(k=1,2,\cdots\) 直到 \(n=\log_2 N\) 执行算式

\[
a_i^{(k)}=a_i^{(k-1)}a_{i-2^{k-1}}^{(k-1)},\quad i=2^k,2^{k+1},\cdots,N-1;
\tag{11.3.10}
\]

\[
b_i^{(k)}=
\begin{cases}
b_i^{(k-1)},& i=0,1,\cdots,2^{k-1}-1,\\
b_i^{(k-1)}+a_i^{(k-1)}b_{i-2^{k-1}}^{(k-1)},& i=2^{k-1},2^{k-1}+1,\cdots,N-1,
\end{cases}
\tag{11.3.11}
\]

\begin{quote}
校注（agent 补充，CH11-E003）：式（11.3.10）中的下标序列确实印为 \(i=2^k,2^{k+1},\cdots,N-1\)，已按放大原图照录。与本页上方推导及式（11.3.7）的逐个下标范围比较，此处第二项疑应为 \(2^k+1\)，而不是 \(2^{k+1}\)；属于原书排印疑误。
\end{quote}

\href{../../source-images/pdf-288.jpeg}{原页图像}

则 \(x_i=b_i^{(n)}\)，\(i=0,1,\cdots,N-1\) 即为所求。

\subsection{本节覆盖原页的校注记录}\label{ux672cux8282ux8986ux76d6ux539fux9875ux7684ux6821ux6ce8ux8bb0ux5f55}

下列记录按原页关联，含同页相邻小节的校注；计算时同时读取上文校注和完整原页。

\begin{itemize}
\tightlist
\item
  \texttt{CH11-E003}：原书 PDF 页 287
\end{itemize}

\end{document}
